An essential requirement for relating logical corpora is standardizing the identifiers so that each identifier in the corpus can be uniquely referenced.
It is desirable to use a uniform naming schema so that the syntax and semantics of identifiers can be understood and implemented as generically as possible.
\mmt URIs have been specifically designed for that purpose, and \autoref{ch:import} shows by example how to systematically generate URIs during importing a library. Additionally, in \cite{MueGauKal:alignmentsreport17} we present URI schemas for a list of selected theorem provers. 

Using these URIs with abbreviated proof assistant
names, the following presents alignments across proof assistants for two representative
concepts.  Also included are some alignments to programming languages, which are relevant
for code generation. Thousands of other alignments, both perfect and imperfect, can be
explored on mathhub~\cite{Alignments:on}.
%In all cases, we will see how big the differences between the details are across the logical corpora even though most of the alignments are in fact perfect.

\paragraph{Cartesian Product}

In constructive type theory, there are two common ways of expressing the non-dependent Cartesian product.
First, if the foundation has inductive types such as the Calculus of Inductive Constructions, it can be an inductive type with one binary constructor. This is the case for:
\begin{compactitem}
\item \texttt{Coq ? Init/Datatypes ? prod.ind}
\item \texttt{Matita ? datatypes/constructors ? Prod.ind}
\end{compactitem}
Second, if the foundation defines a dependent sum type, it is possible to express the Cartesian product as its non-dependent special case:
\begin{compactitem}
\item \texttt{Isabelle ? CTT/CTT ? times}
\end{compactitem}

In higher-order logic, the only way to introduce types is by using the \texttt{typedef} construction, which constructs a new type that is isomorphic to a certain subtype of an existing type.
In particular, most HOL-based systems introduce the Cartesian product $A\times B$ by using a unary predicate on $A\to B\to\cn{bool}$:
\begin{compactitem}
\item \texttt{HOLLight ? pair/type ? prod}
\item \texttt{HOL4 ? pair/type ? prod}
\item \texttt{Isabelle ? HOL/Product ? prod}
\end{compactitem}

In set theory, it is also possible to restrict dependent sum types to obtain the Cartesian product. This approach is used in Isabelle/ZF:

\begin{compactitem}
\item \texttt{Isabelle ? ZF/ZF ? cart\_prod}
\end{compactitem}

In Mizar the Cartesian product is defined as functor in first-order logic. The definition invloves discharging the well-definedness condition. We defined functor is:
\begin{compactitem}
\item \texttt{Mizar ? ZFMISC\_1 ? K2}
\end{compactitem}

In PVS, the product type constructor is part of the system foundation:
\begin{compactitem}
\item \texttt{PVS ? foundation.PVS ? tuple\_tp}
\end{compactitem}

Cartesian products appear also in most programming languages and the code generators
of various proof assistants do use a number of these:
\begin{compactitem}
\item \texttt{OCaml ? core ? * }
\item \texttt{Haskell ? core ? , }
\item \texttt{Scala ? core ? , }
\item \texttt{CPP ? std ? pair }
\end{compactitem}

Informal sources that can be aligned are e.g.:
\begin{compactitem}
\item \small \url{https://en.wikipedia.org/wiki/Cartesian_product}
\item \small \url{https://en.wikipedia.org/wiki/Product_type}
\item \footnotesize \url{http://mathworld.wolfram.com/CartesianProduct.html}
\end{compactitem}

% \paragraph{Addition}

% In constructive type theory, addition of natural numbers is typically defined as a fixed points of certain equations.
% %This is the case for:
% \begin{compactitem}
%  \item \url{Coq ? Init/Nat ? add}
%  \item \url{Matita ? nat/plus ? plus}
% \end{compactitem}

% Higher-order logic proof assistants usually use high-level constructions for primitive
% recursion.  Interestingly, here the alignment is very obvious at this high-level even
% though the elaboration into core language features may result in very different-looking,
% but logically equivalent definitions.  This is how addition is implemented in HOL Light
% and HOL4.  Isabelle/HOL uses a similar construction but inside a type class for
% commutative monoids with difference.
% \begin{compactitem}
%   \item \url{HOLLight ? arith/const ? +}
%   \item \url{HOL4 ? arithmetic/const ? +}
%   \item \url{Isabelle ? HOL/Nat ? plus_nat_inst.plus_nat}
% \end{compactitem}

% In set theory, defining addition is surprisingly the least straightforward: In Isabelle/ZF, addition uses the \texttt{primrec} construction together with a coercion from non naturals to naturals.
% PVS defines a general addition on number fields, which all the concrete number spaces (reals, integers, etc.) inherit.
% In Mizar, it is the restriction of complex addition to natural numbers.
% Mizar's complex addition itself is built from the real number addition, which in turn is built from positive real addition, rational addition, and ordinal addition.
% So in principle it would be possible to align with Mizar's ordinal addition.
% %, however as the ordinal addition is rarely used, we stick to
% %the former:
% \begin{compactitem}
%   \item \url{PVS?Prelude.number_fields?+}
%   \item \url{Isabelle ? ZF/Arith ? add}
%   \item \url{Mizar ? NAT_1 ? K2}%\\
% %    (redefined from \url{Mizar ? XCOMPLEX_0 ? K2})
%  \end{compactitem}

% Most programming languages do not implement arbitrary precision natural
% numbers. However, it is possible to find partial mappings, which are
% in fact already used by efficient code generation~\cite{HaftmannN10} for
% Isabelle/HOL natural numbers and the representation of Coq natural numbers
% by its extended bytecode machine~\cite{ArmandGST10}.
% \begin{compactitem}
% \item \url{OCaml ? Big_Int ? add_big_int }
% \item \url{Haskell ? core ? + }
% \item \url{http://www.smlnj.org/ ? IntInf ? + }
% \end{compactitem}

\paragraph{Concatenation of Lists}

In constructive type theory (e.g. for Matita, Coq), the append operation on lists can
be defined as a fixed point.  In higher-order logic, append for polymorphic lists
can be defined by primitive recursion, as done by HOL Light and HOL4.  Isabelle/HOL
slightly differs from these two because it uses lists that were built with the co-datatype
package~\cite{BlanchetteHLP-ITP14}.
PVS and Isabelle/ZF also use primitive recursion for monomorphic lists.
In Mizar, lists are represented by finite sequences, which are functions from a finite subset of natural numbers (one-based \texttt{FINSEQ} and zero-based \texttt{XFINSEQ}) with append provided.
Concatenation of lists is also common in programming languages.

\begin{compactitem}
%\item Matita: \texttt{cic:/matita/list/list/append}
\item \texttt{Coq ? Init/Datatypes ? app}
\item \texttt{HOLLight ? lists/const ? APPEND}
\item \texttt{HOL4 ? list/const ? APPEND}
\item \texttt{Isabelle ? HOL/List ? append}
%being one of the definitions (special case of
%append for T-sequences).
%\ednote{@CK: this is unclear CK: clarified, is it ok now?}
%\footnote{URIs are omitted for space reasons, they are available online (see Section \ref{sec:system})}
\item \texttt{PVS?Prelude.list\_props?append}
\item \texttt{Isabelle ? ZF/List\_ZF ? app}
\item \texttt{Mizar ? ORDINAL4 / K1} %\\
%  (redefined by \texttt{Mizar ? AFINSQ\_1 / def\_4})
\item \texttt{OCaml ? core ? @ }
\item \texttt{Haskell ? core ? ++ }
\item \texttt{Scala ? core ? ++ }
\end{compactitem}






%%% Local Variables:
%%% mode: latex
%%% TeX-master: "paper-cpp"
%%% End:


%  LocalWords:  ednote compactitem texttt cn tuple_tp Mizar well-definedness monoids
%  LocalWords:  arith plus_nat_inst.plus_nat primrec hspace
